\documentclass{article}

\usepackage{amsmath,amssymb}
\usepackage{xurl}
\usepackage[hidelinks]{hyperref}

% Shared commands for notes in docs/paper-gaps/.
% Notation follows the LDT paper (references/ldt-paper/).

\newcommand{\Lean}{\textsc{Lean}}
\newcommand{\leanid}[1]{\textsf{#1}}
\newcommand{\ghissue}[1]{\href{https://github.com/LionSR/MIPStarRE/issues/#1}{GitHub issue \##1}}
\newcommand{\ghpr}[1]{\href{https://github.com/LionSR/MIPStarRE/pull/#1}{GitHub PR \##1}}

% --- Paper-compatible notation ---
\newcommand{\F}{\mathbb{F}}
\newcommand{\N}{\mathbb{N}}
\newcommand{\C}{\mathbb{C}}
\newcommand{\tr}{\operatorname{tr}}
\newcommand{\ot}{\otimes}
\newcommand{\E}{\mathop{\bf E\/}}
\newcommand{\bx}{\mathbf{x}}
\newcommand{\by}{\mathbf{y}}
\newcommand{\bu}{\mathbf{u}}
\newcommand{\bv}{\mathbf{v}}

% Approximation relation (paper uses \approx_\eps and \simeq_\zeta)
\newcommand{\approxeps}{\approx_{\varepsilon}}
\newcommand{\simgeq}{\simeq_{\zeta}}

% Polynomial spaces (paper uses \polyfunc, \polysub, \polymeas)
\newcommand{\polyfunc}[3]{\operatorname{Poly}(#1,#2,#3)}
\newcommand{\polysub}[3]{\operatorname{PolySub}(#1,#2,#3)}
\newcommand{\polymeas}[3]{\operatorname{PolyMeas}(#1,#2,#3)}


\title{Polynomials as algebraic objects versus functions over finite fields}
\author{}
\date{\today}

\begin{document}
\maketitle

\section{The assertion}

This note concerns the definition of low-individual-degree polynomials
used throughout the low individual degree test
\cite{Ji2020LowIndividualDegree}. The paper defines
$\polyfunc{m}{q}{d}$ as the set of polynomials in $\F_q^m$ with
individual degree at most $d$ (\texttt{preliminaries.tex},
lines~89--94). The Lean formalization has two related concepts:
\begin{itemize}
  \item \texttt{Polynomial params}, a structure wrapping a multivariate
    polynomial (in \texttt{MvPolynomial}) together with a proof that each
    variable's degree is at most $d$;\footnote{\texttt{MIPStarRE/LDT/Basic/LowDegreePolynomial.lean},
      line~13.}
  \item \texttt{HasLowIndividualDegree params g}, a predicate on a function
    $g : \Point \to \Fq$ asserting the existence of a polynomial
    representative with individual degree at most $d$.\footnote{\texttt{MIPStarRE/LDT/Basic/LinePolynomials.lean}, line~20.}
\end{itemize}
The measurement outcome space \texttt{IdxPolyFamily} uses
\texttt{Polynomial params} as the outcome type, i.e., the measurement
outcomes are indexed by algebraic polynomials, not by functions.

\section{The point at issue}

Over a finite field $\F_q$, distinct polynomials can induce the same
function $\F_q^m \to \F_q$. For example, $x^q$ and $x$ are equal as
functions on $\F_q$, but are distinct as algebraic polynomials. If the
outcome space of a measurement were defined as the set of \emph{functions}
of low individual degree, then two outcome labels that represent the same
function would be forced to merge. If it is defined as the set of
\emph{polynomials}, then two distinct polynomial representatives of the
same function are kept separate as measurement outcomes.

The Schwartz--Zippel lemma in the paper (Corollary~2 of
\Cref{lem:schwartz-zippel-total-degree}) states: for distinct
$g, h \in \polyfunc{m}{q}{d}$,
\[
  \Pr_{\bu \sim \F_q^m}[g(\bu) = h(\bu)] \le \frac{md}{q}.
\]
This bound applies to distinct \emph{polynomials} (algebraic objects),
not just distinct functions; it follows from the fact that the difference
$g - h$ is a nonzero polynomial of total degree at most $md$, to which
the Schwartz--Zippel lemma for polynomials is applied.

The Lean formalization uses \texttt{Polynomial params}, whose underlying
type is \texttt{MvPolynomial (Fin params.m) (Scalar params)}. Two
distinct \texttt{MvPolynomial} values that induce the same function on
$\F_q^m$ are stored as distinct outcome labels. The
Schwartz--Zippel bound still applies because the collision probability
of two distinct polynomial objects at a random point is bounded by
$md/q$, exactly as in the paper.

The \texttt{HasLowIndividualDegree} predicate is defined on functions,
but it is used only as a classification tool, not as the outcome space.
Where the paper says ``$g \in \polyfunc{m}{q}{d}$'', the Lean code either
has a \texttt{Polynomial params} value (for indexing measurement outcomes)
or a \texttt{HasLowIndividualDegree} hypothesis (for asserting that a
function comes from a low-degree polynomial). The two notions are linked
by the lemma \texttt{Polynomial.hasLowIndividualDegree}, which extracts
the function from the stored polynomial.

\section{The blueprint discrepancy (resolved)}

A preliminary version of the blueprint defined $\polyfunc{m}{q}{d}$ as a
set of functions $\F_q^m \to \F_q$ represented by polynomials. The
faithfulness audit of 2026-04-05 identified this as a real mathematical
shift, because it would quotient by functional equivalence.\footnote{\texttt{audits/2026-04-05\_faithfulness-audit.md}, lines~33--45.}
The blueprint has since been corrected to read:\footnote{\texttt{blueprint/src/chapter/ch03\_preliminaries.tex}, line~77.}
``we identify elements of $\polyfunc{m}{q}{d}$ with their polynomial
representatives, not their functional equivalence classes.''

The Lean formalization itself never adopted the function-quotient view.
The outcome space has always been algebraic polynomials, matching the
paper's intended definition and the corrected blueprint.

\section{Conclusion}

The Lean formalization uses algebraic polynomials as the measurement
outcome space, which is faithful to the paper. The
\texttt{HasLowIndividualDegree} predicate on functions is a convenience
for stating that a point-evaluated outcome came from a low-degree
polynomial, and does not change the underlying object. The
Schwartz--Zippel applications are sound because they are applied to
distinct polynomial objects, not to distinct functions.

\noindent\textbf{Verdict:} no discrepancy, after the blueprint correction.
The earlier blueprint formulation was a notational misstatement that
has been resolved; the Lean formalization has always been faithful on
this point.

\bibliographystyle{alpha}
\bibliography{docs/paper-gaps/references}

\end{document}